\documentclass[10pt,a4paper]{amsart}
\usepackage[margin=19mm]{geometry}
\usepackage{amsmath,amssymb,mathtools}
\usepackage[scr=rsfs,cal=euler]{mathalfa}
\usepackage{xfrac}
\usepackage[unicode,hidelinks,bookmarks=false]{hyperref}
\newcommand{\ie}{i.e.\ }
\newcommand{\cf}{cf.\ }
\newcommand{\resp}{respectively\ }
\newcommand{\Subsection}{\subsection}
\providecommand{\nobreakdash}{\nobreak\hbox{-}}
\setlength{\emergencystretch}{2em}
\multlinegap0pt
\usepackage{amssymb}
\newtheorem{theorem}{Theorem}
\newtheorem{corollary}{Corollary}
\DeclareMathOperator{\Po}{Po}
\title{P\'olya-Ostrowski Groups and Unit Indices in Real Biquadratic Fields}
\author{Huda Naeem Hleeb Al-Jabbari, Abbas Maarefparvar}
\date{Source: arXiv:2108.01904v4 (CC BY 4.0)}
\begin{document}
\maketitle

\noindent\emph{Source and licence.} P\'olya-Ostrowski Groups and Unit Indices in Real Biquadratic Fields. Version arXiv:2108.01904v4 (CC BY 4.0), distributed by arXiv under the Creative Commons Attribution 4.0 International licence, http://creativecommons.org/licenses/by/4.0/, as recorded in the arXiv licence metadata for this exact version and shown on the abstract page https://arxiv.org/abs/2108.01904v4. Full licence notice: \texttt{licenses/arxiv-2108.01904-CC-BY-4.0.txt}.\par\smallskip

\noindent\textit{Editorial scope.} Counted unit: the article's corollary fields (source0.tex lines 495--506). The article's complete original proof of it is delivered with it (source0.tex lines 508--526, one \texttt{proof} environment of 19 lines). No mathematics is written, repaired or re-derived: the statement and the proof are the article's own lines, in the article's own notation, and no retained line is substituted or re-flowed.  Retained uncounted as the source-local context the proof needs: lines 396--418.  Nothing was omitted from any retained span, so the package carries no omission record.  No retained line contains a citation, so the wrapper carries no bibliography and no bibliography entry.  No retained line contains a figure environment, an image-inclusion command or drawing code, so no image is delivered and the package declares no asset dependency.  Editorial formatting only, in addition: an AMS \texttt{amsart} preamble that restates the article's own declarations and macros used by the retained lines, namely the environments \texttt{theorem}, \texttt{corollary} on separate counters, together with the standard packages those lines need.  The retained environments print in this document as Theorem 1, Corollary 1 in order of appearance, whereas the article prints the same items with the numbering of its own full document.  The complete native source of the article as distributed in the arXiv e-print for this exact version is bundled unchanged as \texttt{sources/114624/source0.tex}.
\begin{theorem} \label{theorem, relatione between order of Polya group and unit index}
Let $K=\mathbb{Q}(\sqrt{d_1},\sqrt{d_2})$ be a real biquadratic field with quadratic subfields $k_i=\mathbb{Q}(\sqrt{d_i})$ ($d_3$ is the squarefree part of $d_1 d_2$). Denote by $t$ the number of quadratic subfields of $K$ whose fundamental unit have negative norm. 
\begin{itemize}
\item[(i)] If $2$ ramifies totally in $K/\mathbb{Q}$ and all $k_i$'s contain elements with the same norm either $2$ or $-2$, then:
\begin{equation*}
    \# \Po(K) = \begin{cases}
               \frac{	\left[U_K : U_{k_1} U_{k_2} U_{k_3}\right]. \prod_{p \mid d_K}e_p(K/\mathbb{Q})}{2^{6-t}}              & : \,  t = 0,1, \\
                & \\ 
               \frac{	\left[U_K : U_{k_1} U_{k_2} U_{k_3}\right]. \prod_{p \mid d_K}e_p(K/\mathbb{Q})}{2^4}              & : \, t =2,3.\\
           \end{cases}
\end{equation*}

\item[(ii)] Otherwise:
\begin{equation*}
    \# \Po(K) = \begin{cases}
               \frac{	\left[U_K : U_{k_1} U_{k_2} U_{k_3}\right]. \prod_{p \mid d_K}e_p(K/\mathbb{Q})}{2^{5-t}}              & : \,  t = 0,1, \\
               & \\
               \frac{	\left[U_K : U_{k_1} U_{k_2} U_{k_3}\right]. \prod_{p \mid d_K}e_p(K/\mathbb{Q})}{2^3}              & : \, t=2,3.\\
           \end{cases}
\end{equation*}
\end{itemize}
In particular, for $t\in \{2,3\}$  P\'olya groups have the same order.
\end{theorem}

\begin{corollary} \label{corollary, improving Zantema upper bound for Polya biquadratic fields}
	Let
$K$ be a real biquadratic field. Denot by $t$ the number of quadratic subfields of $K$ whose fundamental units have negative norm, and by $s_K$ the number of primes that are ramified in $K/\mathbb{Q}$. If $K$ is P\'olya, then
\begin{equation*}
s_K \leq \left\{
\begin{array}{rl}
5 \, :& \, t=0, \\
 4 \, :& \,   t=1, \\
 3 \, :& \,  t=2,3.
\end{array} \right.
\end{equation*}
\end{corollary}

\begin{proof}
 Suppose that $\Po(K)=0$. We prove the cases $t=2,3$. The other two cases, say $t=0,1$, can be proved similarly. Let $k_1,k_2,k_3$ be the three quadratic subfield of $K$, and suppose that $k_1$ and $k_2$ have some units of negative norm. If $2$ is totally ramified in $K/\mathbb{Q}$, and all $k_i$'s contain elements with the same norm either $2$ or $-2$, then by part (i) of Theorem \ref{theorem, relatione between order of Polya group and unit index}, we get
\begin{equation*}
\prod_{p | d_K} e_p(K/\mathbb{Q})=e_2(K/\mathbb{Q}) \cdot \prod_{\substack{p | d_K\\ p \neq 2}} e_p(K/\mathbb{Q}) = \frac{2^4}{\left[U_K : U_{k_1} U_{k_2} U_{k_3}\right]},
\end{equation*}
which implies that 
\begin{equation*}
 \prod_{\substack{p | d_K\\ p \neq 2}} e_p(K/\mathbb{Q}) \leq 2^2,
\end{equation*}
or equivalently, $s_K \leq 3$. On the other hand, if $2$ doesn't ramify totally in $K/\mathbb{Q}$ or one of the equations
\begin{equation*}
N_{k_i/\mathbb{Q}}(x_i)=+2 \quad \text{or} \quad -2, \quad i=1,2,3,
\end{equation*}
doesn't have any solution $x_i \in \mathcal{O}_{k_i}$, then by part (ii) of Theorem \ref{theorem, relatione between order of Polya group and unit index}, we obtain
\begin{equation*}
	\prod_{p | d_K} e_p(K/\mathbb{Q}) = \frac{2^3}{\left[U_K : U_{k_1} U_{k_2} U_{k_3}\right]} \leq 2^3,
\end{equation*}
and again we get $s_K \leq 3$.
\end{proof}

\end{document}
